\chapter{Conclusion et perspectives}
\label{chap:conclusion}

\section{Bilan du travail}

Ce stage avait pour objet de transformer une décision quotidienne prise au jugement en une
décision calculée, reproductible et démontrablement optimale. Les quatre objectifs fixés au
chapitre~\ref{chap:problematique} ont été atteints.

\paragraph{Formaliser.} Le procédé industriel a été traduit en un programme linéaire mixte
en nombres entiers comprenant six ensembles d'indices, dix-sept familles de contraintes et
douze hypothèses explicitement tracées. Chaque contrainte est reliée à un organe physique
identifié, ce qui rend le modèle discutable avec un exploitant --- condition nécessaire pour
qu'il soit un jour utilisé.

\paragraph{Résoudre.} Le modèle atteint l'optimum global en \SI{0,06}{\second} et satisfait
\SI{100}{\percent} de la demande, largement en deçà de la contrainte de temps compatible avec
un usage quotidien.

\paragraph{Valider.} Un validateur écrit indépendamment du modèle, à partir des seules
équations physiques, confirme la solution par quatre-vingt-six contrôles. Ce choix
architectural n'était pas théorique : il a détecté un défaut réel dès la première exécution.

\paragraph{Restituer.} Les rapports Excel et JSON sont produits au format imposé par le
service, et le programme est utilisable dans une chaîne automatisée.

\section{Apports méthodologiques}

Trois résultats dépassent le cadre strict du livrable attendu.

\subsection{L'identification de l'unité de compte}

Le fait que toutes les grandeurs du système soient exprimées en tonnes de \PdeuxOcinq{}
n'est écrit dans aucun des neuf documents fournis. Cette information a dû être
\emph{déduite}, par trois preuves indépendantes et convergentes. Sans elle, l'intégralité
des bilans matière serait fausse et le modèle produirait des plans irréalisables.

\begin{encadre}
La leçon dépasse ce projet : \emph{les conventions les plus structurantes d'un système sont
souvent celles que personne ne pense à écrire}, parce qu'elles vont de soi pour ceux qui y
travaillent. Les débusquer fait partie du travail de modélisation.
\end{encadre}

\subsection{La démonstration de dégénérescence de l'objectif}

L'objectif spécifié --- maximiser le total livré sous contrainte de satisfaction exacte de
la demande --- est constant sur le domaine réalisable. Sa dérivée est nulle ; il ne sélectionne
aucune solution.

Ce défaut n'était pas anodin : il aurait produit un outil renvoyant des plans instables d'un
jour à l'autre sans raison apparente, et donc rapidement abandonné par ses utilisateurs. La
correction proposée --- un objectif lexicographique à trois niveaux --- généralise l'objectif
demandé plutôt que de le remplacer, ce qui permet de la défendre devant le prescripteur.

\subsection{L'audit des spécifications}

Onze anomalies ont été identifiées, dont quatre erreurs arithmétiques atteignant
\SI{59}{\percent} d'écart. Deux enseignements en découlent, applicables à tout projet :
\emph{recalculer plutôt que recopier}, et \emph{fixer une hiérarchie de sources avant
d'examiner les cas}, pour éviter tout arbitrage de circonstance.

\section{Le résultat le plus utile pour l'exploitation}

Paradoxalement, la contribution principale de ce travail n'est pas le plan optimal
quotidien, mais le \emph{diagnostic} qu'il a permis d'établir.

Le modèle sert toute la demande sans difficulté. Mais il révèle deux violations
structurelles --- le débordement d'IR11 de \SI{1016}{\tonne} par jour et le sous-remplissage
chronique du stock de 14EXT --- qu'\emph{aucune} décision d'exploitation ne peut corriger,
parce qu'elles découlent d'un paramètre de configuration et non d'une variable de décision.

L'analyse de sensibilité en établit la cause unique et le remède chiffré : ramener de quatre
à un le nombre d'échelons de 14EXT affectés à la cocristallisation résorbe intégralement les
deux violations, sans investissement et sans dégrader le service.

\begin{encadre}[title={Ce qu'un modèle d'optimisation doit produire}]
Non pas seulement un plan, mais \textbf{le diagnostic de ce qui empêche le plan d'être bon}.
Une solution optimale qui viole les capacités physiques n'est pas un échec du modèle : c'est
une information de premier ordre sur le système réel.
\end{encadre}

\section{Sur la méthode de travail}

Le découpage en sept phases, avec l'interdiction de coder ce qui n'a pas d'abord été
formalisé, a produit un effet mesurable : les phases~3 à~6 --- socle logiciel, modèle,
validation, restitution --- ont été menées en une seule itération de travail, alors que
quatre à six étaient estimées.

L'explication est directe. Le modèle mathématique ayant été posé contrainte par contrainte,
avec pour chacune sa traduction en français et sa justification, l'écriture du code s'est
réduite à une \emph{traduction}. Les erreurs qui coûtent cher --- confusion entre paramètre
et variable, oubli d'un terme de bilan, mauvaise interprétation d'une règle métier --- avaient
déjà été éliminées sur le papier, où elles coûtent peu.

\section{Limites}

\paragraph{L'horizon mono-période.} Le modèle optimise une journée en tenant compte de
l'état final des stocks, mais sans anticiper la demande du lendemain. Un enchaînement de
décisions localement optimales n'est pas nécessairement optimal sur la semaine.

\paragraph{Les paramètres en attente.} Deux données restent inconnues. L'analyse de
sensibilité montre que l'une d'elles est sans influence, mais l'autre --- la part minimale
de DEC\_CL dans IR11 --- pilote directement le niveau de décadmiation.

\paragraph{Le déterminisme.} Rendements, production et demande sont supposés connus avec
certitude. Un aléa de production ou une casse d'équipement invalideraient le plan du jour.

\paragraph{La validation terrain.} Les résultats n'ont pas encore été confrontés à ce que
fait réellement l'usine un jour donné. C'est l'étape indispensable avant tout déploiement.

\section{Perspectives}

\subsection{Rendre la cocristallisation décidable}

C'est la perspective la plus directement rentable. L'analyse de sensibilité montre que
l'affectation des échelons à la cocristallisation --- aujourd'hui un paramètre --- détermine
à elle seule les deux violations structurelles. En faire une \textbf{variable de décision
binaire par échelon} permettrait au modèle de trouver seul la configuration qui les évite.

Le coût est faible : une vingtaine de binaires supplémentaires, négligeable pour le solveur.
La difficulté est ailleurs : il faut d'abord vérifier avec l'exploitation qu'une telle
reconfiguration est opérationnellement possible, et à quelle fréquence.

\subsection{L'extension multi-période}

Étendre l'horizon à une semaine transformerait le modèle en un problème de planification
avec stocks de liaison entre périodes. La formulation actuelle s'y prête : il suffirait
d'indexer les variables par le temps et de remplacer le terme de cible du niveau 3 par un
véritable couplage inter-périodes. La taille du problème resterait modeste.

\subsection{La prise en compte de l'incertitude}

Une approche par scénarios ou une optimisation robuste permettrait de produire des plans
tenant compte de la variabilité des rendements et des aléas de production.

\subsection{L'industrialisation}

Le passage d'un outil d'étude à un outil d'exploitation supposerait une interface de saisie
des scénarios, une connexion aux systèmes de mesure des niveaux de cuve, et un historique
des plans produits permettant de comparer le réalisé au planifié.

\section{Apport personnel}

Ce stage m'a confronté à une réalité que la formation académique aborde peu : \emph{un cahier
des charges réel est incomplet et parfois faux}. Le réflexe d'accepter les données fournies
telles quelles aurait conduit à un modèle silencieusement erroné --- ce qui est bien plus
dangereux qu'un modèle visiblement cassé.

J'en retiens trois pratiques que je compte conserver : recalculer systématiquement plutôt
que recopier ; ne jamais se vérifier avec le même raisonnement que celui qui a produit le
résultat ; et transformer une incertitude en analyse de sensibilité plutôt qu'en hypothèse
dissimulée.

Sur le fond, ce travail m'a permis de mesurer la distance entre un modèle correct et un
modèle \emph{utile}. Un modèle correct satisfait ses contraintes. Un modèle utile explique à
son utilisateur ce qui le limite, et lui propose une action.
